\section{Empirical Discovery: Mean-Dominated Geometry}
\label{sec:empirical}

This section presents the central empirical finding of the paper:
activation tensors in large language models contain a dense, coherent mean
component that dominates extreme element-wise outliers and materially
contributes to spectral anisotropy.
We contrast this behavior with gradient-output tensors and establish that
bias is activation-specific.

\subsection{First-Order Structure of Training Signals}
\label{subsec:first_order_asymmetry}

Training of large language models is driven by two primary data objects at each layer:
the forward activations $x_i \in \mathbb{R}^H$ and the backward gradient outputs
$\delta_i \in \mathbb{R}^H$.
These two tensors jointly determine the parameter update through outer-product
structure
\[
g_i = \delta_i x_i^\top.
\]
Any systematic structure present in either $x_i$ or $\delta_i$
directly propagates into the gradient matrix and therefore into the model parameters.

We begin by examining their first-order statistics.
For a batch of $N$ tokens, define the empirical mean
\[
\mu_z = \frac{1}{N} \sum_{i=1}^N z_i,
\qquad z_i \in \{x_i, \delta_i\}.
\]

To determine whether $\mu_z$ reflects coherent structure or merely finite-sample
fluctuation, we compare its magnitude to the expected size of a noise-level mean.
If $\{z_i\}$ were independent zero-mean samples with variance
$\mathbb{E}\|z_i\|_2^2$, then by standard concentration,
\[
\|\mu_z\|_2 \approx \sqrt{\frac{\mathbb{E}\|z_i\|_2^2}{N}}.
\]
This motivates the noise-normalized mean magnitude
\[
R_z
=
\frac{\|\mu_z\|_2}
{\sqrt{\mathbb{E}\|z_i\|_2^2 / N}}.
\]

Interpretation:
\begin{itemize}
\item $R_z \approx 1$ indicates that the empirical mean is consistent with
finite-sample noise.
\item $R_z \gg 1$ indicates coherent, batch-aligned structure.
\end{itemize}

\paragraph{Empirical measurement.}

Across layers (3 and 21) and training steps (10k and 180k),
we observe a pronounced asymmetry:

\begin{itemize}
\item Activations exhibit large coherent means,
      with $R_x$ ranging from approximately $17$ to $41$.
\item Gradient outputs remain near noise level,
      with $R_\delta$ ranging from approximately $1.05$ to $1.19$.
\end{itemize}

Figure~\ref{fig:activation_grad_R} reports these values.

\begin{figure}[t]
\centering
\includegraphics[width=0.7\linewidth]{figures/sec2_1_R_bar.png}
\caption{
Noise-normalized mean magnitude $R_z$ for activations and gradient outputs.
Activations exhibit strong coherent first-order structure ($R_x \gg 1$),
whereas gradient outputs remain consistent with noise ($R_\delta \approx 1$).
}
\label{fig:activation_grad_R}
\end{figure}

\paragraph{Observation.}

The data reveal a clear first-order asymmetry:
coherent mean structure is present in activations but absent in gradient outputs.
Since gradients are constructed as outer products $g_i = \delta_i x_i^\top$,
this implies that coherent first-order bias in training must originate
from the activation side.

This asymmetry motivates the subsequent analysis.
If coherent bias dominates the forward signal but not the backward signal,
then anisotropy observed in gradients and optimizer statistics
must be driven primarily by activation structure.

\subsection{Outliers Are Dominated by the Activation Mean}
\label{subsec:outlier_mean_dominance}

We now identify the origin of extreme elementwise magnitudes in training tensors.
While spectral analyses describe variance structure at the matrix level,
low-bit quantization and several optimizer pathologies are triggered by
\emph{entrywise} extremes.
We therefore ask:

\begin{quote}
\emph{Among the largest-magnitude entries, which component contributes most:
mean bias, spike variance, or long-tail variance?}
\end{quote}

\paragraph{Decomposition.}
For each matrix $X \in \mathbb{R}^{N\times H}$ (activation or gradient output),
we decompose
\[
X = M + X_{\text{spike}} + X_{\text{tail}},
\]
where $M=\mathbf{1}\mu^\top$ is the broadcasted mean component,
$X_{\text{spike}}$ is the projection of the centered matrix onto the top-$k$
right singular vectors (we use $k=1\%$ of $H$),
and $X_{\text{tail}}$ is the remaining component.

\paragraph{Outlier set and contribution shares.}
We define outliers as the top $0.1\%$ entries ranked by raw magnitude $|X_{ij}|$.
For each outlier entry, we compute the squared contribution shares
\[
\rho_{\text{mean}} = \frac{M_{ij}^2}{X_{ij}^2},\qquad
\rho_{\text{spike}} = \frac{(X_{\text{spike}})_{ij}^2}{X_{ij}^2},\qquad
\rho_{\text{tail}} = \frac{(X_{\text{tail}})_{ij}^2}{X_{ij}^2}.
\]
The resulting distributions of $\rho_{\text{mean}},\rho_{\text{spike}},\rho_{\text{tail}}$
provide a direct elementwise attribution of extreme magnitudes.

\paragraph{Results across layers and training steps.}
Figure~\ref{fig:outlier_contrib_grid} shows the outlier contribution histograms
for activations (left column) and gradient outputs (right column),
across two layers (3 and 21) and two training steps (10k and 180k).

Two consistent patterns emerge:
\begin{itemize}
    \item \textbf{Activations: mean-dominated outliers.}
    Activation outliers concentrate strongly near $\rho_{\text{mean}}\approx 1$,
    indicating that the mean component accounts for the overwhelming majority of
    squared magnitude in the most extreme entries.
    Spike contributions are secondary and tail contributions are negligible.
    \item \textbf{Gradient outputs: variance-dominated outliers.}
    Gradient-output outliers show negligible mean contribution
    ($\rho_{\text{mean}}\approx 0$ for nearly all outliers) and are instead
    dominated by variance components (spike and/or tail).
\end{itemize}

\paragraph{Implication.}
This establishes the central empirical fact behind this paper:
\begin{quote}
\emph{Elementwise outliers in activations are overwhelmingly bias-driven, whereas
gradient-output outliers are variance-driven.}
\end{quote}
Consequently, mean subtraction targets the primary source of quantization-relevant
activation outliers, while downstream spectral conditioning addresses residual
variance anisotropy.

\begin{figure*}[t]
\centering

% ================= Row 1 =================
\begin{subfigure}[t]{0.23\linewidth}
  \centering
  \includegraphics[width=\linewidth]{figures/layer3_step010000_activation.pt_outlier_hist.png}
  \caption{L3 10k (Act)}
\end{subfigure}
\hfill
\begin{subfigure}[t]{0.23\linewidth}
  \centering
  \includegraphics[width=\linewidth]{figures/layer3_step010000_grad_output.pt_outlier_hist.png}
  \caption{L3 10k (Grad)}
\end{subfigure}
\hfill
\begin{subfigure}[t]{0.23\linewidth}
  \centering
  \includegraphics[width=\linewidth]{figures/layer3_step180000_activation.pt_outlier_hist.png}
  \caption{L3 180k (Act)}
\end{subfigure}
\hfill
\begin{subfigure}[t]{0.23\linewidth}
  \centering
  \includegraphics[width=\linewidth]{figures/layer3_step180000_grad_output.pt_outlier_hist.png}
  \caption{L3 180k (Grad)}
\end{subfigure}

\vspace{0.6em}

% ================= Row 2 =================
\begin{subfigure}[t]{0.23\linewidth}
  \centering
  \includegraphics[width=\linewidth]{figures/layer21_step010000_activation.pt_outlier_hist.png}
  \caption{L21 10k (Act)}
\end{subfigure}
\hfill
\begin{subfigure}[t]{0.23\linewidth}
  \centering
  \includegraphics[width=\linewidth]{figures/layer21_step010000_grad_output.pt_outlier_hist.png}
  \caption{L21 10k (Grad)}
\end{subfigure}
\hfill
\begin{subfigure}[t]{0.23\linewidth}
  \centering
  \includegraphics[width=\linewidth]{figures/layer21_step180000_activation.pt_outlier_hist.png}
  \caption{L21 180k (Act)}
\end{subfigure}
\hfill
\begin{subfigure}[t]{0.23\linewidth}
  \centering
  \includegraphics[width=\linewidth]{figures/layer21_step180000_grad_output.pt_outlier_hist.png}
  \caption{L21 180k (Grad)}
\end{subfigure}

\caption{
Outlier contribution decomposition across layers and training steps.
Each panel shows squared contribution shares for the top $0.1\%$ largest entries.
Across all settings, activation outliers are overwhelmingly mean-dominated,
whereas gradient-output outliers are variance-dominated.
}
\label{fig:outlier_contrib_grid}

\end{figure*}


\subsection{Spectral Role of Bias}
\label{subsec:spectral_role_bias}

Having established that activation bias dominates extreme entries, we now
investigate how this bias manifests in the singular value spectrum.

Let $X \in \mathbb{R}^{N \times H}$ denote the activation matrix.
We compute the singular value spectrum of both:

\begin{itemize}
\item the raw activation matrix $X$,
\item the mean-centered activation matrix $X_c = X - \mathbf{1}\mu^\top$.
\end{itemize}

Our goal is to determine:
\begin{enumerate}
\item how much of the spectral spike is attributable to bias,
\item whether bias accounts for the entire spike,
\item how this behavior varies across layers and training stages.
\end{enumerate}

\paragraph{Experimental setup.}
For each of four configurations
(layer 3 and layer 21 at 10k and 180k steps),
we compute:

\begin{itemize}
\item singular value spectrum (log-log scale),
\item cumulative spectral energy.
\end{itemize}

The spike rank $k$ is defined as $1\%$ of the feature dimension.

\paragraph{Results.}

\begin{figure}[t]
\centering
\begin{subfigure}[t]{0.24\linewidth}
\includegraphics[width=\linewidth]{figures/layer3_10k_spectrum_raw_vs_centered.png}
\caption{Layer 3, 10k (spectrum)}
\end{subfigure}
\begin{subfigure}[t]{0.24\linewidth}
\includegraphics[width=\linewidth]{figures/layer3_10k_cum_energy_raw_vs_centered.png}
\caption{Layer 3, 10k (energy)}
\end{subfigure}
\begin{subfigure}[t]{0.24\linewidth}
\includegraphics[width=\linewidth]{figures/layer3_180k_spectrum_raw_vs_centered.png}
\caption{Layer 3, 180k (spectrum)}
\end{subfigure}
\begin{subfigure}[t]{0.24\linewidth}
\includegraphics[width=\linewidth]{figures/layer3_180k_cum_energy_raw_vs_centered.png}
\caption{Layer 3, 180k (energy)}
\end{subfigure}

\vspace{0.5em}

\begin{subfigure}[t]{0.24\linewidth}
\includegraphics[width=\linewidth]{figures/layer21_10k_spectrum_raw_vs_centered.png}
\caption{Layer 21, 10k (spectrum)}
\end{subfigure}
\begin{subfigure}[t]{0.24\linewidth}
\includegraphics[width=\linewidth]{figures/layer21_10k_cum_energy_raw_vs_centered.png}
\caption{Layer 21, 10k (energy)}
\end{subfigure}
\begin{subfigure}[t]{0.24\linewidth}
\includegraphics[width=\linewidth]{figures/layer21_180k_spectrum_raw_vs_centered.png}
\caption{Layer 21, 180k (spectrum)}
\end{subfigure}
\begin{subfigure}[t]{0.24\linewidth}
\includegraphics[width=\linewidth]{figures/layer21_180k_cum_energy_raw_vs_centered.png}
\caption{Layer 21, 180k (energy)}
\end{subfigure}

\caption{
Singular spectra and cumulative energy before and after mean removal.
Mean removal significantly suppresses the leading singular component
in early layers, while leaving the long tail largely unchanged.
}
\label{fig:spectral_bias_role}
\end{figure}

\paragraph{Observation 1: Bias strongly contributes to spike energy in early layers.}

In layer 3 at 10k steps, mean removal dramatically reduces the leading singular
value and shifts the cumulative energy curve downward for small $k$.
This confirms that a substantial portion of the spectral spike is due to
the coherent mean component.

\paragraph{Observation 2: Bias is not the entire spike.}

Even after mean removal, a clear low-rank structure remains.
The spike does not vanish entirely.
This indicates that anisotropy contains two components:

\begin{itemize}
\item a dense coherent mean component,
\item structured low-rank variance.
\end{itemize}

\paragraph{Observation 3: Bias influence weakens in deeper layers.}

For layer 21 at 10k, raw and centered spectra are nearly identical.
This indicates that deep-layer representations contain substantially less
mean-driven anisotropy.

\paragraph{Observation 4: Bias evolves during training.}

Comparing 10k and 180k steps shows that bias-driven spike strength
changes during training.
Early stages exhibit stronger mean-dominated structure.

\paragraph{Conclusion.}

The spectral spike decomposes into:

\begin{equation}
\text{Spike} = \text{Dense coherent mean} + \text{Low-rank structured variance}.
\end{equation}

Mean removal eliminates the first term but leaves the second intact.
Therefore:

\begin{itemize}
\item Bias is a structured component inside the spike,
\item but it is not synonymous with the entire spike,
\item and spike variance persists after mean removal.
\end{itemize}

This refined understanding will be essential in Section~\ref{sec:grad_propagation},
where we analyze how bias propagates into gradient updates.

\subsection{Bias Propagation into Gradient Spectral Geometry}
\label{subsec:grad_spectrum_bias}

Having established that activation tensors contain a statistically significant
dense mean component (Section~2.1) and that this mean dominates extreme
activation outliers (Section~2.2), we now examine how this bias propagates into
the gradient matrix that directly drives parameter updates.

\paragraph{Gradient structure.}
For a linear layer, the per-token gradient has outer-product structure
\[
g_i = \delta_i x_i^\top,
\qquad
G = \frac{1}{N}\sum_i g_i = \frac{1}{N}\Delta^\top X.
\]
Although gradient outputs $\delta_i$ exhibit no significant mean bias
(Section~2.1), the activation matrix $X$ does.
Therefore, even if $\Delta$ is approximately zero-mean, the product
$\Delta^\top X$ may inherit structured anisotropy through $X$.

To isolate this effect, we compare:
\[
G_{\text{raw}} = \frac{1}{N}\Delta^\top X,
\qquad
G_{\text{centered}} = \frac{1}{N}\Delta^\top (X - \mu_x).
\]

\paragraph{Experimental protocol.}
For each layer and training step, we compute:
\begin{itemize}
\item the singular value spectrum of $G_{\text{raw}}$,
\item the singular value spectrum of $G_{\text{centered}}$,
\item the cumulative spectral energy distribution.
\end{itemize}
All spectra are plotted on log–log scale to reveal spike–tail separation.

\begin{figure}[t]
\centering

% Row 1
\begin{subfigure}[t]{0.24\linewidth}
\includegraphics[width=\linewidth]{figures/layer3_10k_G_spectrum_raw_vs_centered.png}
\caption{Layer 3, 10k (spectrum)}
\end{subfigure}
\begin{subfigure}[t]{0.24\linewidth}
\includegraphics[width=\linewidth]{figures/layer3_10k_G_cum_energy_raw_vs_centered.png}
\caption{Layer 3, 10k (cumulative)}
\end{subfigure}
\begin{subfigure}[t]{0.24\linewidth}
\includegraphics[width=\linewidth]{figures/layer3_180k_G_spectrum_raw_vs_centered.png}
\caption{Layer 3, 180k (spectrum)}
\end{subfigure}
\begin{subfigure}[t]{0.24\linewidth}
\includegraphics[width=\linewidth]{figures/layer3_180k_G_cum_energy_raw_vs_centered.png}
\caption{Layer 3, 180k (cumulative)}
\end{subfigure}

\vspace{0.5em}

% Row 2
\begin{subfigure}[t]{0.24\linewidth}
\includegraphics[width=\linewidth]{figures/layer21_10k_G_spectrum_raw_vs_centered.png}
\caption{Layer 21, 10k (spectrum)}
\end{subfigure}
\begin{subfigure}[t]{0.24\linewidth}
\includegraphics[width=\linewidth]{figures/layer21_10k_G_cum_energy_raw_vs_centered.png}
\caption{Layer 21, 10k (cumulative)}
\end{subfigure}
\begin{subfigure}[t]{0.24\linewidth}
\includegraphics[width=\linewidth]{figures/layer21_180k_G_spectrum_raw_vs_centered.png}
\caption{Layer 21, 180k (spectrum)}
\end{subfigure}
\begin{subfigure}[t]{0.24\linewidth}
\includegraphics[width=\linewidth]{figures/layer21_180k_G_cum_energy_raw_vs_centered.png}
\caption{Layer 21, 180k (cumulative)}
\end{subfigure}

\caption{Gradient singular spectra before and after activation mean removal.
Across layers and training stages, removing the activation mean consistently
reduces the dominance of the leading singular directions and shifts spectral
energy toward the long tail.}
\label{fig:grad_spectrum_bias}
\end{figure}

\paragraph{Observations.}

Across all layers and training stages, three consistent phenomena emerge:

\begin{enumerate}
\item \textbf{Raw gradients exhibit spike–tail structure.}
Despite the absence of mean bias in gradient outputs $\delta_i$,
the gradient matrix $G$ displays strong anisotropy, confirming that
activation geometry propagates into the optimizer.

\item \textbf{Mean removal weakens spike dominance.}
After centering activations, the leading singular values shrink,
and cumulative energy curves shift rightward.
The spike becomes less dominant relative to the tail.

\item \textbf{The spike does not vanish.}
Even after mean removal, a structured low-rank component remains.
Thus, the spectral spike consists of two parts:
\[
\text{Spike} = \text{Dense mean component} + \text{Structured covariance component}.
\]
Mean removal eliminates the dense bias term but leaves the structured
covariance component intact.
\end{enumerate}

\paragraph{Interpretation.}

These results establish a key structural asymmetry:

\begin{itemize}
\item Activation bias directly inflates gradient anisotropy.
\item Removing activation bias reduces gradient spike energy
without affecting gradient-output statistics.
\item The remaining spike reflects genuine covariance structure rather
than trivial mean effects.
\end{itemize}

This separation is crucial for understanding optimization behavior:
activation bias acts as a low-rank energy amplifier that artificially
inflates second-moment statistics, while structured covariance encodes
learning-relevant variability.

In Section~2.4B we quantify how this bias inflation impacts adaptive
second-moment optimizers.

\subsection{Parameter Matrices Do Not Accumulate Bias}
\label{subsec:parameter_bias}

Previous sections
established that activation tensors exhibit strong coherent mean bias,
and that this bias propagates multiplicatively into gradient matrices.
We now examine whether such bias accumulates in the model parameters
themselves during training.

\paragraph{Measurement protocol.}
For representative feed-forward layers (Layer~3 and Layer~21),
we analyze weight matrices
$W \in \mathbb{R}^{3072 \times 1024}$
at early (10k) and later (180k) training steps.
We compute both column-wise and row-wise mean components and evaluate:

\begin{enumerate}
\item Noise-normalized mean magnitude
\[
R = \frac{\|\mu\|_2}{\sqrt{\mathbb{E}\|W\|_F^2 / d}},
\]
where $d$ denotes the corresponding dimensionality.
\item Mean energy fraction
\[
\rho = \frac{\|\mathbf{1}\mu^\top\|_F^2}{\|W\|_F^2}.
\]
\end{enumerate}

\begin{table}[t]
\centering
\caption{Bias statistics of model parameters across layers and training stages.}
\label{tab:weight_bias}
\begin{tabular}{lcccc}
\toprule
Layer & Step & $R_{\text{col}}$ & $\rho_{\text{col}}$ & $\rho_{\text{row}}$ \\
\midrule
Layer 3  & 10k   & 0.970 & $3.06\times10^{-4}$ & $9.85\times10^{-4}$ \\
Layer 3  & 180k  & 0.988 & $3.18\times10^{-4}$ & $9.89\times10^{-4}$ \\
Layer 21 & 10k   & 0.999 & $3.25\times10^{-4}$ & $9.81\times10^{-4}$ \\
Layer 21 & 180k  & 1.002 & $3.27\times10^{-4}$ & $9.65\times10^{-4}$ \\
\bottomrule
\end{tabular}
\end{table}

\paragraph{Observations.}

Across all layers and training stages:

\begin{itemize}
\item Noise-normalized mean ratios remain at $R \approx 1$,
indicating that parameter means are indistinguishable from
finite-sample noise.
\item Mean energy fractions remain below $0.1\%$ of total Frobenius energy.
\item No systematic growth of parameter bias is observed over training.
\end{itemize}

In contrast, activation tensors exhibit
$R_x \in [17, 41]$ (Section~\ref{subsec:mean_asymmetry}),
demonstrating that coherent bias is specific to intermediate
representations rather than learned parameters.

\paragraph{Implication.}
These results establish that mean-dominated geometry is not
a persistent property of model weights.
Bias originates in activations, influences gradients transiently,
but does not accumulate into parameter matrices.
This distinction localizes the geometric source of bias-driven
anisotropy to the forward representation space.

\subsection{Bias Originates at the Embedding Layer}
\label{subsec:embedding_bias}

To determine whether activation bias is induced by deeper layers or
already present at the input stage, we analyze the embedding-layer
activations, i.e., the token representations fed into the transformer
stack.

Let $x_{b,t} \in \mathbb{R}^H$ denote the embedding of token $t$ in
sequence $b$, and let $X \in \mathbb{R}^{(B L)\times H}$ be the
flattened batch matrix.
We measure bias magnitude using the same noise-normalized statistic
introduced in Section~2.1:
\[
R \;=\;
\frac{\|\mu\|_2}
{\sqrt{\mathbb{E}\|x\|_2^2 / N}},
\qquad
\mu = \frac{1}{N}\sum_{i=1}^N x_i,
\]
where $N = B L$.

\paragraph{Batch-level bias.}
For the embedding layer we observe:
\[
R_{\text{batch}} \approx 24.0 \text{ (step 10k)},
\qquad
R_{\text{batch}} \approx 25.6 \text{ (step 180k)}.
\]
Since $R\approx 1$ corresponds to noise-level mean,
these values indicate that embedding activations possess a
highly significant coherent mean component.

\paragraph{Per-sequence bias.}
We further compute sequence-level means
\[
\mu_b = \frac{1}{L} \sum_{t=1}^L x_{b,t}
\]
for each of the $B=8$ sequences in the batch.

The average per-sequence noise-normalized ratio is:
\[
\mathbb{E}_b[R_{\text{seq}}] \approx 8.8 \text{ (10k)},
\qquad
\mathbb{E}_b[R_{\text{seq}}] \approx 9.3 \text{ (180k)}.
\]

Thus, bias is already strong within individual sequences,
and is not merely a batching artifact.

\paragraph{Cross-sequence coherence.}
To quantify directional consistency, we measure the cosine
similarity between each sequence mean and the global batch mean:
\[
\cos(\mu_b, \mu_{\text{batch}}).
\]

We obtain:
\[
\text{mean cosine} \approx 0.968 \text{ (10k)},
\qquad
\text{mean cosine} \approx 0.976 \text{ (180k)}.
\]

These near-unity values indicate that sequence-level biases
are almost perfectly aligned across sequences.

\paragraph{Implications.}
These results establish three structural facts:

\begin{enumerate}
\item The embedding layer already contains a strong coherent bias.
\item This bias is present within individual sequences.
\item The bias direction is globally aligned across sequences.
\end{enumerate}

Therefore, activation bias is not created by deeper nonlinear
transformations; rather, it originates at the embedding stage
and propagates through the network.

This observation is critical:
the geometric anisotropy observed in deeper layers is seeded
by a coherent input bias rather than emerging purely from
layerwise covariance structure.

\subsubsection{Causal Origin of Activation Bias: Token Frequency Structure}
\label{subsec:token_frequency_bias}

The preceding subsections established that activation bias is
(i) statistically significant,
(ii) responsible for extreme outliers,
and (iii) absent in gradient-output tensors and learned weight matrices.
We now identify its structural origin.

\paragraph{Frequency-weighted embedding mean.}

Let $E \in \mathbb{R}^{V \times H}$ denote the learned embedding table,
and let $p(v)$ denote the empirical token frequency in a batch.
The activation mean induced by token embeddings is

\[
\mu_{\text{freq}}
=
\sum_{v=1}^{V} p(v)\, E[v].
\]

This expression shows that activation bias is not arbitrary:
it is the frequency-weighted average of embedding vectors.
If high-frequency tokens share coherent geometric structure,
$\mu_{\text{freq}}$ becomes large.

\paragraph{Concentration of bias contribution.}

Using $1{,}024 \times 1{,}024$ sampled tokens
($1{,}048{,}576$ total tokens),
we measured how much of $\|\mu_{\text{freq}}\|_2$
is explained by the top-$K$ most frequent tokens.

Results show:

\begin{itemize}
\item Top 10 tokens explain $\approx 62\%$ of $\|\mu_{\text{freq}}\|$.
\item Top 20 tokens explain $\approx 83\%$.
\item Top 50 tokens explain $\approx 110\%$ (due to constructive alignment).
\item Top 200 tokens explain $\approx 132\%$.
\end{itemize}

This demonstrates that activation bias is highly concentrated
in a small subset of frequent tokens.

\paragraph{Identity of dominant tokens.}

Table~\ref{tab:top_tokens_bias} lists the top contributing tokens
by absolute bias contribution.
These tokens are overwhelmingly function words and structural markers.

\begin{table}[t]
\centering
\caption{Top tokens contributing to embedding-induced activation bias (Layer 3, 180k step).}
\label{tab:top_tokens_bias}
\begin{tabular}{rll}
\toprule
Rank & Token ID & Decoded Token \\
\midrule
1  & 151643 & \texttt{<|endoftext|>} \\
2  & 279    & \texttt{the} \\
3  & 11     & \texttt{,} \\
4  & 13     & \texttt{to} \\
5  & 311    & \texttt{and} \\
6  & 323    & \texttt{of} \\
7  & 315    & \texttt{a} \\
8  & 264    & \texttt{in} \\
9  & 220    & \texttt{is} \\
10 & ...    & \texttt{that} \\
\midrule
\multicolumn{3}{c}{Additional high contributors: \texttt{you, for, it, I, on, are, with, be, as, 's, have, was, or, (}} \\
\bottomrule
\end{tabular}
\end{table}

These tokens correspond to:

\begin{itemize}
\item Articles and determiners
\item Prepositions
\item Conjunctions
\item Auxiliary verbs
\item Punctuation
\item Special sequence markers
\end{itemize}

In other words, activation bias aligns strongly with grammatical structure.

\paragraph{Geometric interpretation.}

The bias vector can therefore be interpreted as a
\emph{grammar direction} in embedding space:

\[
\mu_{\text{freq}} \approx
\sum_{v \in \text{function tokens}} p(v)\, E[v].
\]

Since grammar tokens appear in nearly every sentence,
their contributions accumulate coherently across sequences,
producing a dense rank-one component in activation space.

\paragraph{Connection to earlier findings.}

This explains the empirical observations:

\begin{enumerate}
\item Activation bias magnitude ratios $R_x \gg 1$ (Section~2.1).
\item Mean-dominated outliers (Section~2.2).
\item Significant spectral energy reduction after centering (Section~2.3).
\item Gradient bias inflation via outer-product propagation (Section~2.4).
\end{enumerate}

Crucially, the bias does \emph{not} arise from
optimizer dynamics or weight accumulation;
it originates directly from corpus frequency imbalance
interacting with coherent embedding geometry.

\paragraph{Conclusion.}

Activation bias is a deterministic consequence of:

\[
\text{Zipfian token frequency}
\;\times\;
\text{coherent embedding structure}.
\]

This provides a complete causal explanation for
mean-dominated outliers in large language model training.

\subsection{Mean Removal as the First-Order Mechanism for Preserving Long-Tail Resolution}
\label{sec:mean_first_order_quantization}

Low-bit quantization allocates limited dynamic range across all activation values.
Under blockwise quantization, the quantization scale is determined by extreme
entries within each block.
Consequently, the component responsible for extreme magnitudes directly
controls representational resolution of the remaining features.

This section provides quantitative evidence that activation mean bias,
rather than residual spike variance, is the dominant driver of
long-tail degradation under low-bit quantization.

\subsubsection{Experimental Setup}

We evaluate four activation tensors:

\begin{itemize}
\item Layer 3 at step 10k (early layer, early training),
\item Layer 3 at step 180k (early layer, late training),
\item Layer 21 at step 10k (deep layer, early training),
\item Layer 21 at step 180k (deep layer, late training).
\end{itemize}

For each tensor, we:

\begin{enumerate}
\item Decompose activation into mean, spike (top 1\%), and tail components.
\item Identify the top 0.1\% largest elementwise magnitudes.
\item Compute squared contribution shares of mean, spike, and tail.
\item Perform blockwise INT4 quantization (block size = 64).
\item Measure tail signal-to-quantization-noise ratio (SQNR) at $p99.9$.
\end{enumerate}

\subsubsection{Outlier Decomposition}

Across all four settings, we observe:

\begin{itemize}
\item The mean component explains approximately $85\%$–$90\%$
of squared magnitude among the top 0.1\% largest activation entries.
\item The spike component explains roughly $5\%$–$15\%$.
\item The tail contribution is negligible.
\end{itemize}

Thus, extreme values that determine quantization scale
are overwhelmingly mean-dominated.

\subsubsection{Quantitative Impact on Tail SQNR}

We compare three conditions:

\begin{itemize}
\item Raw activation.
\item Mean removed.
\item Spike removed (mean preserved).
\end{itemize}

Representative tail SQNR results (block=64, $p99.9$):

\begin{center}
\begin{tabular}{lccc}
\toprule
Setting & Raw & Mean Removed & Spike Removed \\
\midrule
Layer 3 (10k)   & 8.7 dB  & 12.0 dB  & 10.8 dB \\
Layer 3 (180k)  & 10.6 dB & 11.9 dB  & 11.0 dB \\
Layer 21 (10k)  & 11.2 dB & 13.0 dB  & 12.0 dB \\
Layer 21 (180k) & 11.1 dB & 13.5 dB  & 12.0 dB \\
\bottomrule
\end{tabular}
\end{center}

Two consistent quantitative patterns emerge:

\begin{enumerate}
\item Removing the mean improves tail SQNR by approximately
$+2$ to $+4$ dB across settings.
\item Removing spike components while keeping the mean yields
smaller improvements.
\end{enumerate}

Therefore, the mean bias is the primary determinant
of long-tail quantization degradation.

\subsubsection{First-Order vs Second-Order Effects}

Although spike components may explain a nontrivial portion
of total spectral energy,
they contribute far less to extreme elementwise magnitudes
that determine quantization scale.

Formally:
\[
\text{Outlier Energy} \approx
\text{Mean Energy} \gg \text{Spike Energy} \gg \text{Tail Energy}.
\]

Since block quantization scale is governed by
\[
\alpha \propto \max_{i \in \text{block}} |x_i|,
\]
and these maxima are mean-dominated,
mean removal directly reduces dynamic-range inflation.

Spike whitening provides secondary refinement,
but does not address the principal source of extreme values.

\subsubsection{Conclusion}

These results establish that:

\begin{quote}
\textbf{Activation mean bias is the first-order mechanism
responsible for long-tail resolution loss under low-bit quantization.}
\end{quote}

Removing the mean alone yields the largest improvement
in tail SQNR with minimal computational cost.
Residual spike treatment is beneficial but secondary.

This finding motivates addressing bias at its source
(embedding activation)
rather than relying on repeated layer-wise spectral decompositions.

\subsection{Activation Bias and Second-Moment Inflation in Adaptive Optimization}
\label{subsec:optimizer_bias}

We now quantify how activation bias affects adaptive optimizers.
While gradient outputs exhibit no statistically significant mean
(Section~2.1), the activation tensor does.
Because per-token gradients are outer products
\[
g_i = \delta_i x_i^\top,
\]
any dense mean component in $x_i$ directly propagates into gradient
statistics.

\paragraph{Adam second-moment structure.}
Adam and AdamW maintain an elementwise second-moment estimate
\[
v \approx \mathbb{E}[g_i \odot g_i].
\]
Using the identity
\[
v = \frac{1}{N}(\Delta^{\odot 2})^\top (X^{\odot 2}),
\]
we see that squared activation magnitudes directly influence the
scale of $v$.

If $x = \mu_x + \tilde{x}$, then
\[
g_i = \delta_i \mu_x^\top + \delta_i \tilde{x}_i^\top,
\]
and the first term produces a dense, coherent contribution to $v$.

\paragraph{Experimental protocol.}
For each layer and training stage we compute:
\begin{itemize}
\item $v_{\text{raw}}$ from original activations,
\item $v_{\text{mean}}$ after mean removal,
\item $v_{\text{mean+whiten}}$ after mean removal and top-$1\%$ spike whitening.
\end{itemize}
We report extreme-tail statistics ($p99.9$ and $\max$), which dominate
learning-rate suppression in Adam.

\begin{table}[t]
\centering
\caption{Extreme-tail statistics of Adam second moment $v$.}
\label{tab:adam_v_bias}
\begin{tabular}{lcccc}
\toprule
Layer & Step & Method & $p99.9(v)$ & $\max(v)$ \\
\midrule
Layer 3 & 10k & Raw & $1.60\times10^{-10}$ & $8.38\times10^{-10}$ \\
 & & Mean-only & $3.44\times10^{-11}$ & $1.91\times10^{-10}$ \\
 & & Mean+whiten & $1.72\times10^{-11}$ & $8.09\times10^{-11}$ \\
\midrule
Layer 3 & 180k & Raw & $1.00\times10^{-10}$ & $7.37\times10^{-9}$ \\
 & & Mean-only & $2.50\times10^{-11}$ & $1.53\times10^{-10}$ \\
 & & Mean+whiten & $2.44\times10^{-11}$ & $1.54\times10^{-10}$ \\
\midrule
Layer 21 & 10k & Raw & $1.67\times10^{-11}$ & $3.43\times10^{-11}$ \\
 & & Mean-only & $1.59\times10^{-11}$ & $2.58\times10^{-11}$ \\
 & & Mean+whiten & $1.48\times10^{-11}$ & $2.33\times10^{-11}$ \\
\midrule
Layer 21 & 180k & Raw & $9.23\times10^{-11}$ & $1.10\times10^{-9}$ \\
 & & Mean-only & $2.21\times10^{-11}$ & $2.69\times10^{-10}$ \\
 & & Mean+whiten & $2.09\times10^{-11}$ & $1.95\times10^{-10}$ \\
\bottomrule
\end{tabular}
\end{table}

\paragraph{Key observations.}

\begin{enumerate}
\item \textbf{Mean removal dramatically reduces extreme $v$ outliers.}
In early layers and early training (Layer 3, 10k), the $p99.9$
statistic decreases by more than $4\times$, and the maximum by more
than $4\times$.
In some later-stage cases, the reduction in $\max(v)$ approaches
$50\times$.

\item \textbf{Median values change minimally.}
Across all runs, the median of $v$ changes by less than 1\%.
This indicates that bias primarily inflates a small subset of
extreme coordinates rather than the bulk of parameters.

\item \textbf{Whitening provides additional but smaller gains.}
In early training stages, spike whitening further reduces extreme
values.
In later stages, most of the pathological inflation is already
removed by mean subtraction.
\end{enumerate}

\paragraph{Interpretation.}

These results demonstrate that activation bias is a dominant driver
of extreme second-moment inflation in adaptive optimizers.
Even though gradient outputs themselves are nearly zero-mean,
the dense activation bias propagates multiplicatively through
outer-product gradients and produces heavy-tailed $v$ statistics.

Removing the activation mean significantly stabilizes second-moment
estimates and therefore mitigates excessive effective learning-rate
suppression in Adam-style optimizers.

This provides a direct mechanistic link between activation bias
and optimization pathologies.


\subsection{Bias-Induced Dynamic Range Distortion in Low-Bit Quantization}
\label{subsec:quant_bias}

Low-bit quantization is fundamentally sensitive to dynamic range.
Under blockwise quantization, each block of size $B$ uses a shared
scale parameter proportional to
\[
\alpha = \max_{j \in \text{block}} |x_j|.
\]
If a small subset of entries exhibits unusually large magnitude,
they determine $\alpha$, thereby compressing the effective resolution
allocated to the majority of entries.

Given the decomposition
\[
x = \mu_x + x_{\text{spike}} + x_{\text{tail}},
\]
we now empirically evaluate how much of the dynamic range distortion
is attributable to activation bias.

\paragraph{Experimental protocol.}
For each layer and training stage we compute:

\begin{itemize}
\item Entrywise magnitude statistics ($p99$, $p99.9$, $\max$),
\item Blockwise $p99.9$ of block-maximum magnitudes (block size 64),
\item Reduction ratios after mean removal,
\item Additional reductions after spike whitening.
\end{itemize}

\begin{table}[t]
\centering
\caption{Dynamic range statistics before and after bias removal.}
\label{tab:quant_bias}
\begin{tabular}{lcccc}
\toprule
Layer & Step & Method & Entry $p99.9$ & Block $p99.9$ \\
\midrule
Layer 3 & 10k & Raw & $6.24$ & $10.40$ \\
 & & Mean-only & $3.09$ & $5.27$ \\
 & & Mean+whiten & $2.34$ & $3.22$ \\
\midrule
Layer 3 & 180k & Raw & $2.29$ & $3.42$ \\
 & & Mean-only & $1.16$ & $1.71$ \\
 & & Mean+whiten & $1.06$ & $1.61$ \\
\midrule
Layer 21 & 10k & Raw & $3.52$ & $4.60$ \\
 & & Mean-only & $3.46$ & $4.52$ \\
 & & Mean+whiten & $4.00$ & $4.00$ \\
\midrule
Layer 21 & 180k & Raw & $4.17$ & $6.48$ \\
 & & Mean-only & $2.31$ & $4.20$ \\
 & & Mean+whiten & $2.98$ & $2.98$ \\
\bottomrule
\end{tabular}
\end{table}

\paragraph{Key observations.}

\begin{enumerate}
\item \textbf{Mean removal consistently reduces extreme magnitudes.}
In early-layer early-training regimes (Layer 3, 10k),
entrywise $p99.9$ reduces by approximately $2\times$,
and blockwise $p99.9$ reduces by approximately $2\times$.
This indicates that a substantial portion of dynamic range inflation
originates from the mean component.

\item \textbf{Later layers exhibit smaller bias sensitivity.}
Layer 21 at 10k shows minimal reduction after mean removal,
consistent with the weaker bias ratio observed in Section~2.1.

\item \textbf{Spike whitening further smooths residual anisotropy.}
In some regimes, spike whitening provides additional reduction
beyond mean removal, but the primary distortion mechanism remains
bias-dominated.
\end{enumerate}

\paragraph{Interpretation.}

These results demonstrate that activation bias directly inflates
the blockwise dynamic range that governs low-bit quantization.
Removing the bias reduces extreme magnitudes by approximately
$2\times$ in sensitive regimes, thereby restoring effective
resolution for long-tail semantic components.

Importantly, the majority of dynamic range distortion is attributable
to the mean component rather than the long-tail spectrum itself.

This provides a concrete mechanistic explanation for why mean removal
substantially improves training stability in FP4 and related
low-bit settings.